% Раздел: Используемые данные (глава «Метод и экспериментальная установка»)

\subsection{Используемые данные}
\label{sec:datasets}

B-LoRA относится к однофотографическим подходам: для файнтюна достаточно одного изображения-референса. Экспериментальная оценка дополнительно требует референсного распределения для FID и набора контентных промптов для семантической оценки. Сводка наборов приведена в таблице~\ref{tab:datasets}.

\textbf{Эталонные стили (DS1).} Отобрано 8~изображений двух художественных школ~--- импрессионизма Моне и постимпрессионизма Ван~Гога,~--- получены с WikiArt и выбраны по критериям стилистической выраженности и цветовой характерности. Для Ван~Гога: «Звёздная ночь» (1889), «Красные виноградники» (1888), «Подсолнухи» (1888), «Автопортрет с перевязанным ухом» (1889); для Моне: «Стог сена в~Живерни» (1886), «Водяные лилии~--- вечерний эффект» (1899), «Мост Ватерлоо» (1903), «Мадам Моне с ребёнком» (ок.~1875). При каждом запуске обучения используется одно изображение; набор из четырёх позволяет оценить вариативность результатов внутри стиля. Выбор обусловлен стилистической контрастностью: импрессионизм Моне отличается размытыми контурами и пастельной палитрой, манера Ван~Гога~--- вихреобразной фактурой и насыщенным цветом.

\begin{figure}[h]
    \centering
    \includegraphics[width=0.95\textwidth]{style_refs_grid.png}
    \caption{Стилевые референсы. Верхний ряд~--- Ван~Гог: «Красные виноградники в~Арле» (1888), «Автопортрет с~перевязанным ухом» (1889), «Подсолнухи» (1888), «Звёздная ночь» (1889). Нижний ряд~--- Моне: «Стог сена в~Живерни» (1886), «Мадам Моне с~ребёнком» (ок.~1875), «Водяные лилии~--- вечерний эффект» (1899), «Мост Ватерлоо» (1903). Источник: WikiArt; работы в общественном достоянии.}
    \label{fig:style_refs}
\end{figure}

\textbf{ArtBench-10 (DS3).} Датасет~\cite{liao2022artbench} (около 60\,000~изображений по 10~стилям) используется для оценки FID в \textit{reduced protocol}: подвыборка по 500~изображений на жанр. Из-за ограниченного объёма метрика применяется для сравнения конфигураций между собой, а не для сопоставления с опубликованными значениями. Источник: HuggingFace Hub (\texttt{Doub7e/ArtBench-10}), CC-BY.

\textbf{MS-COCO val2017 (DS2).} Подвыборка из 100~промптов валидационного раздела~\cite{lin2014coco} обеспечивает нейтральный контентный фон для измерения CLIP-сходств. Источник: \url{cocodataset.org}, CC-BY.

\textbf{DreamBooth (DS5).} Репозиторий \texttt{google/dreambooth}~\cite{ruiz2023dreambooth}, Apache~2.0. Используются 8~предметов: \texttt{dog} ($\approx$30~изображений~--- основной объект Phase~1b) и 7~дополнительных (\texttt{cat}, \texttt{backpack}, \texttt{bowl}, \texttt{can}, \texttt{clock}, \texttt{vase}, \texttt{bear}, по~$\approx$25~изображений). Совокупный пул~--- $\approx$200~контентных изображений для количественного сравнения методов (Phase~5.2).

\textbf{Парный набор DS8.} Воспроизводит протокол раздела~5.1 оригинальной статьи~\cite{frenkel2024blora} в виде \textit{reduced protocol}: $N_c{=}8$ объектов из DS5, $N_s{=}8$ стилей из DS1; из декартова произведения $8{\times}8{=}64$ семплируется 50~пар при \texttt{seed=42}. Манифест \texttt{experiments/data/b\_lora\_eval\_pairs.json} версионирован через DVC. Основная метрика~--- среднее косинусное сходство DINO~ViT-B/8 между сгенерированным изображением и стилевым референсом. Поскольку пул отличается от оригинального, абсолютные значения не сопоставимы напрямую с опубликованной Table~1; относительные различия между методами остаются корректными.

\begin{table}[h]
    \centering
    \caption{Сводка использованных данных}
    \label{tab:datasets}
    \begin{tabular}{p{2.8cm}p{4.2cm}p{4.0cm}p{3.4cm}}
        \hline
        \textbf{Датасет} & \textbf{Назначение} & \textbf{Размер} & \textbf{Лицензия} \\
        \hline
        WikiArt (DS1) & Обучение B-LoRA; DS8 стили & 8~изображений (2~стиля~$\times$~4) & общественное достояние \\
        DreamBooth~\cite{ruiz2023dreambooth} (DS5) & DS8 контент; Phase~1b & 8~предметов ($\approx$200~изображений) & Apache 2.0 \\
        ArtBench-10~\cite{liao2022artbench} (DS3) & Референс FID (approximate) & 500~изображений на жанр (reduced) & CC-BY \\
        MS-COCO val2017~\cite{lin2014coco} (DS2) & Инференс; CLIP-content & 100~промптов (подвыборка) & CC-BY \\
        DS8: парный eval & Сравнение методов (Phase~5.2) & 8~$\times$~8~$\to$~50~пар (reduced) & DS1~+~DS5 \\
        \hline
    \end{tabular}
\end{table}

Все наборы версионированы средствами DVC: \texttt{.dvc}-файлы фиксируют хэши содержимого в git-репозитории, а сами данные хранятся в S3-совместимом удалённом хранилище. Команда \texttt{dvc pull} восстанавливает идентичный набор после клонирования.
